% !TeX root = ../../../main.tex
\section{A Reproducible Execution of the Audit Procedure}\label{app:demo}

This appendix reports one complete, deliberately small execution of Algorithm~1 (\cref{sec:algorithm}) and one quantitative rejection case, produced by \texttt{scripts/demo\_tfim\_audit.py} in the paper's artifact; \texttt{scripts/verify\_calculations.py} independently reproduces every number quoted in \cref{sec:resources}. The quantum-execution portion runs through \emph{Ariadion}, an \textbf{author-developed, version-pinned research software artifact} (commit \texttt{e4c56fb16b35382c95ba54d4b2b9f1fd2c684683}, version 0.1.0rc2). Because the simulator and the paper share an author, the artifact is \emph{never its own sole validation}: every Ariadion output is cross-checked against a separately implemented pure-NumPy classical oracle that performs exact diagonalization and an independent state-vector simulation of the same circuit, importing no Ariadion code. The same circuits are also rebuilt and executed in Qiskit. The evidence in this appendix falls into three distinct classes that must not be conflated: (i) \emph{classical simulation} of an ideal or declared-toy-noise device, which validates the audit procedure and the shot-allocation mathematics; (ii) independent Qiskit cross-validation, including simulation against a noise model built from a real device's calibration snapshot; and (iii) a single uncorrected execution on one physical device, which is minimal empirical evidence with shot-noise-only uncertainty and must not be read more broadly. None of it is an advantage claim. The instance is chosen so that the honest audit outcome is \textsc{classical-only}---a framework that cannot reject its own demonstration instance would not be worth trusting.

\subsection{Acceptance audit instance: 4-qubit transverse-field Ising model}

\textbf{Gate 1 (contract).} Estimate the energy of a prepared state of the open-chain transverse-field Ising model $H=-J\sum_{i=1}^{3}Z_iZ_{i+1}-h\sum_{i=1}^{4}X_i$ with $J=h=1$, to additive error $\epsilon=0.05$ with joint failure probability $\delta=0.05$. The oracle's exact diagonalization of the $16\times16$ matrix gives ground energy $E_0=-4.758770$, the classical baseline.

\textbf{Gates 2--3 (mechanism, interface).} The prepared state is produced by an explicit fixed circuit---$R_Y(\theta_1)$ on all four qubits, a CX chain, $R_Y(\theta_2)$ on all four, with $\theta_1=1.114887$, $\theta_2=0.206420$ chosen once offline---whose exact energy is $\langle\psi|H|\psi\rangle=-4.653959$, within $0.105$ of $E_0$. The observable decomposes into $G=2$ commuting measurement groups measured by two Ariadion programs: a joint $Z$-basis setting for all $ZZ$ terms (weight $a_Z=3$) and a terminal-Hadamard setting for all $X$ terms (weight $a_X=4$), $\sum_g a_g=7$. Ariadion returns the \emph{joint} classical distribution over the four returned bits per setting; the reduction of those joint distributions to a grouped energy estimate is performed by the paper's artifact code (Ariadion exposes no public expectation-value helper), which retains covariance among observables sharing a setting. A dedicated bit-order probe circuit confirms the least-significant-bit convention assumed by the reduction. The estimator identity $\hat H=\hat\mu_Z+\hat\mu_X$ holds to machine precision, and Ariadion's ideal state-vector distributions agree with the independent oracle to $10^{-9}$ in every outcome probability. Group standard deviations are $\sigma_Z=1.545$, $\sigma_X=1.421$, both well below their worst-case bounds $a_g$.

\textbf{Gate 4 (physics).} The two allocations of \cref{sec:resources} give, at $(\epsilon,\delta)=(0.05,0.05)$:
worst-case \eqref{eq:worstcase-shots}: $N_Z=N_X=171{,}776$, total $343{,}552$ shots;
variance-aware \eqref{eq:neyman-shots} with the exact $\sigma_g$: total $13{,}517$ shots---a $25\times$ reduction, consistent with $\sigma_g<a_g$. Seeded Ariadion sampled runs were executed at every reported shot count; the worst-case-allocation estimate landed within $0.004$ of the exact energy, inside the $\epsilon$ budget. A 400-trial coverage experiment (drawn by the artifact's seeded sampler from Ariadion's exported ideal joint distributions, since 400 repetitions of $3.4\times10^5$-shot runs exceed a pure-Python simulator's practical budget) gave empirical coverage $\Pr[|\hat H-\langle H\rangle|\leq\epsilon]=1.000$ for the worst-case allocation and $0.955$ for the Neyman allocation against the $0.95$ target: the bound over-delivers, the asymptotic approximation lands near its nominal confidence. This is the guarantee distinction of \cref{sec:resources} made empirical. Density-matrix runs under a \emph{declared toy noise model} (depolarizing channels on the one-qubit $R_Y$ and $H$ gates, symmetric $1\%$ readout error) shift the energy from $-4.654$ to $-4.432$ at $p=0.005$ and $-4.129$ at $p=0.02$, monotonically toward zero as expected; Ariadion's density path binds channels to one-qubit gates only, so \textbf{the CX gates execute noiselessly, and the model omits leakage, correlated errors, and any device-calibration ingestion}---it is a toy sensitivity check, not a hardware model, and Ariadion's protection-requirement reporting is deliberately not used because energy estimation has no natural accept/reject event. A toy screen \eqref{eq:nisq-screen} for a plausible hardware ansatz (24 one-qubit gates, 9 two-qubit gates, 4 measurements at representative error rates) gives $P_0=0.887$: the instance would fit comfortably on hardware.

\textbf{Gate 5 (opponent) and decision.} The classical baseline solves the instance exactly in microseconds at every accuracy. No regime exists in which the quantum route wins, and executing the circuits through a real SDK stack changes nothing about that. \emph{Decision: \textsc{classical-only}.} All five gates executed; the procedure produced a defensible rejection with every intermediate quantity checkable against exact values.

\subsection{Rejection case: DNA phase encoding at Gate 3}

The phase-encoding proposal of \cref{sec:audits} was audited quantitatively at its most favorable operating point: a strongly periodic sequence ($N=4096$ bases, period-16 motif), where quantum Fourier transform (QFT) sampling has a sharp peak to find. This case is a purely classical analysis (no Ariadion execution). The simulation confirms the mechanism works as derived---with $10^4$ samples from the exact QFT output distribution, the dominant frequency was correctly identified at a multiple of $N/16=256$. It also confirms the cost structure that kills the proposal: preparing \cref{eq:phase-sequence} from classical data costs $\Theta(N)$ per shot, so the $10^4$ shots imply roughly $4.1\times10^{7}$ loading operations, while a classical fast Fourier transform returns the \emph{complete} spectrum once for about $N\log_2 N\approx4.9\times10^{4}$ operations. The route fails Gate~3 under the stated access model even when Gate~2 succeeds. \emph{Decision: \textsc{classical-only}}, unless the access model changes (for example, data generated coherently rather than loaded).

\subsection{Evidence artifacts and reproducibility}

The run emits a machine-readable evidence set in \texttt{artifacts/}: compiled circuit IR with source-to-operation provenance for both settings; exact ideal distributions and pre-observation amplitudes; seeded sampled counts at every reported shot count; density-matrix distributions, execution provenance, and Ariadion noise-impact reports at both noise strengths; Theonoe state and execution-trace inspections (state-vector path only---trace capture is unsupported on the density path); the independent oracle's results; a manifest recording the resolved Ariadion commit, Python version, full dependency lock (\texttt{pip freeze}), backend, execution modes, every seed, the noise configuration, and the density-path omissions listed above; and SHA-256 hashes of every artifact file. Six numerical tests are asserted in the script: Ariadion-versus-oracle agreement, correct bit ordering, the stated confidence experiment, bit-identical reproduction under identical seeds, preservation of the \textsc{classical-only} decision, and the documented response to a change in modeled noise.

To reproduce from a clean environment, follow the paper repository's README, which documents installing the Ariadion monorepo's sibling packages together at the commit pinned at the head of this appendix and then running \texttt{scripts/demo\_tfim\_audit.py}.

\subsection{Independent cross-validation and the hardware path}

Beyond the Ariadion runs above, three further evidence sets are archived alongside them: an independent SDK cross-validation, a measured-calibration noise simulation, and a single physical-device execution.

\emph{Independent SDK cross-validation.} The same two measurement-setting circuits were rebuilt from scratch in Qiskit (\texttt{scripts/hw\_tfim\_audit.py}) and executed on the Aer simulator. Three-way agreement is asserted programmatically at tolerance $10^{-9}$: Qiskit's ideal statevector distributions against the NumPy oracle, and against Ariadion's archived exact distributions. Ariadion's simulation is therefore validated not only by the oracle but by an unrelated, widely used SDK. Seeded Aer sampling at the Neyman budget reproduced the reference energy within $0.7$ shot-noise standard errors ($\hat{E}=-4.637\pm0.026$).

\emph{Measured-calibration noise.} The circuits were transpiled (optimization level 3) to IBM's Heron-class \texttt{torino} topology and executed against a noise model constructed from that device's actual calibration snapshot---per-gate error rates, $T_1$/$T_2$, and readout errors, all archived verbatim in \texttt{artifacts/hardware/}. The result is consistent with the failure mode the framework describes for uncorrected hardware: $\hat{E}=-4.490\pm0.026$, a $+0.164$ bias toward zero, against a calibration-derived Gate-4 screen of $P_0=0.940$ computed per executed operation from the measured error rates (versus $0.887$ from the representative rates used in \cref{sec:resources}). Counts, transpiled QASM3, calibration snapshots, job identifiers, and SHA-256 manifests are archived per backend.

\emph{Physical-hardware execution.} The identical script's hardware path was run once on IBM's 156-qubit Heron-class \texttt{ibm\_kingston} processor (open plan, least-busy selection) at the full Neyman budget (7{,}042 $Z$-setting and 6{,}475 $X$-setting shots). The result: $\hat{E}=-4.432\pm0.027$ against $E_{\mathrm{ref}}=-4.654$, a $+0.222$ bias toward zero---consistent with the failure signature the audit describes for uncorrected hardware, observed on a physical device with a favorable device-day Gate-4 screen of $P_0=0.971$ computed from \texttt{kingston}'s calibration at run time. The quoted uncertainty is the archived shot-noise standard error only; it excludes calibration drift, temporal correlations, crosstalk, and other systematic effects, which a single uninterleaved run cannot separate from the bias. The noise simulation in the preceding paragraph used a \emph{different} device's calibration (\texttt{torino}), so the two numbers illustrate the same qualitative failure mode but do not constitute a controlled model-versus-device comparison; that experiment---a matched \texttt{kingston}-snapshot simulation against repeated or interleaved \texttt{kingston} blocks, with shot-noise and between-run variation reported separately---is specified but is not part of this paper's evidence. Raw counts, transpiled QASM3, the calibration snapshot, IBM job identifiers, and SHA-256 manifests are archived in \texttt{artifacts/hardware/hardware\_ibm\_kingston/}. \Cref{fig:hardware-energies} plots all three execution contexts against the exact reference.

\begin{figure}[htbp]
\centering
\begin{tikzpicture}
\begin{axis}[
  width=0.88\textwidth,
  height=5.6cm,
  xmin=0.4, xmax=3.6,
  xtick={1,2,3},
  xticklabels={%
    Aer\\ideal sampling,
    calibrated noise\\\texttt{torino} snapshot,
    physical device\\\texttt{ibm\_kingston}},
  x tick label style={font=\footnotesize,align=center,text width=3.1cm},
  ylabel={energy estimate $\hat{E}$},
  ylabel style={font=\small},
  y tick label style={font=\footnotesize},
  ymin=-4.72, ymax=-4.36,
  ytick={-4.7,-4.6,-4.5,-4.4},
  grid=major,
  grid style={draw=softgray},
  axis line style={draw=ink},
  tick style={draw=ink},
  clip=false
]
\addplot[quantumdark,dashed,thick,forget plot] table[x=index,y=eref] {src/figures/hardware-energies.dat};
\node[anchor=west,font=\footnotesize,text=quantumdark] at (axis cs:3.05,-4.654) {$E_{\mathrm{ref}}=-4.654$};
\addplot+[
  only marks,
  mark=*,
  mark size=2.2pt,
  color=quantum,
  error bars/.cd, y dir=both, y explicit,
  error bar style={line width=0.8pt,draw=quantum}
] table[x=index,y=ehat,y error=stderr] {src/figures/hardware-energies.dat};
\end{axis}
\end{tikzpicture}
\caption{Archived energy estimates for the four-qubit demonstration instance across three execution contexts, with the exact reference energy $E_{\mathrm{ref}}$ (dashed). Error bars are the archived shot-noise standard errors only; they exclude calibration drift, temporal correlations, and crosstalk, so the two noisy points must not be read as a controlled model-versus-device comparison---the noise model and the device are different processors. The monotone drift toward zero is the failure signature the audit anticipates for uncorrected hardware, not an advantage result. Every plotted value is regenerated from \texttt{artifacts/hardware/} by \texttt{scripts/make\_figures.py}.}
\label{fig:hardware-energies}
\end{figure}

\emph{Scale behavior.} A sweep of the same audit pipeline over chain lengths $n=4$ to $24$ (\texttt{artifacts/scale/}) records where the classical baselines actually stand. The one-dimensional transverse-field Ising model is exactly solvable by a free-fermion transformation \citep{pfeuty1970}, so the strongest matched classical opponent computes the exact ground energy in $O(n^{3})$ time---milliseconds at every $n$ in the sweep, with no wall anywhere in reach; the sweep implements this baseline and cross-checks it against dense diagonalization at every $n\leq14$. Generic dense diagonalization ($O(8^{n})$ time, $O(4^{n})$ memory), shown for contrast, exhausts practicality near $n=14$ (about two minutes)---a property of the structure-blind method, \emph{not} of the instance family. Statevector evaluation stays within seconds through $n=24$, sampled simulation at 4{,}096 shots stays sub-second, and the measurement budgets remain modest under a constant per-site precision rule: the worst-case Hoeffding allocation \eqref{eq:worstcase-shots} at $\delta=0.05$ and the Neyman variance-target allocation \eqref{eq:neyman-shots} at $z_{1-\delta/2}$ are both computed per $n$ and archived with their formulas in the manifest. Any claim that this family walls classically would be false; a scale claim for the audit pipeline on a genuinely hard family requires a nonintegrable instance benchmarked against tensor-network methods, which this sweep does not attempt. \Cref{fig:scale-sweep} plots the archived timings.

\begin{figure}[htbp]
\centering
\begin{tikzpicture}
\begin{axis}[
  width=0.88\textwidth,
  height=6.4cm,
  xlabel={chain length $n$},
  ylabel={wall-clock time (s)},
  xlabel style={font=\small},
  ylabel style={font=\small},
  tick label style={font=\footnotesize},
  xmin=3, xmax=25,
  xtick={4,8,12,16,20,24},
  ymode=log,
  ymin=1e-5, ymax=1e3,
  grid=major,
  grid style={draw=softgray},
  axis line style={draw=ink},
  tick style={draw=ink},
  legend style={font=\footnotesize,at={(0.02,0.98)},anchor=north west,draw=softgray},
  legend cell align=left,
  unbounded coords=jump
]
\addplot[quantum,thick,mark=*,mark size=1.6pt] table[x=n,y=freefermion] {src/figures/scale-timings.dat};
\addlegendentry{free-fermion baseline, $O(n^{3})$}
\addplot[warning,thick,mark=square*,mark size=1.6pt] table[x=n,y=exactdiag] {src/figures/scale-timings.dat};
\addlegendentry{dense diagonalization, $O(8^{n})$}
\addplot[quantumdark,thick,dashed,mark=triangle*,mark size=2pt] table[x=n,y=statevector] {src/figures/scale-timings.dat};
\addlegendentry{statevector evaluation}
\addplot[ink,thick,dotted,mark=diamond*,mark size=2pt] table[x=n,y=aer4096] {src/figures/scale-timings.dat};
\addlegendentry{sampled simulation, 4{,}096 shots}
\end{axis}
\end{tikzpicture}
\caption{Archived scale sweep. The matched classical opponent---the exact free-fermion solution---stays in the sub-millisecond band across the whole range and shows no wall. Dense diagonalization, plotted only for contrast, leaves the range near $n=14$; that is a property of the structure-blind method, not of the instance family, and it is why a crossover argued against dense diagonalization would be an artifact of choosing a weak baseline. The sweep is therefore evidence \emph{against} a quantum case for this family. Regenerated from \texttt{artifacts/scale/} by \texttt{scripts/make\_figures.py}; the dense-diagonalization series is not computed beyond $n=14$.}
\label{fig:scale-sweep}
\end{figure}

What this appendix still does \emph{not} establish: quantum advantage of any kind (the hardware run demonstrates the audit pipeline and its noise triage on a physical device, not a speedup), a controlled calibration-model-versus-device comparison (the noise simulation and the hardware run used different devices), error-mitigated or error-corrected hardware performance, systematic-uncertainty characterization beyond shot noise, and the framework's value on instances where the decision is genuinely uncertain. That harder demonstration is specified in \cref{sec:audits}.
